\chapter{Probabilidad y variables aleatorias}\label{ch:g11:prob}

\omterm{def:g10:proba:distribution}{probabilidad} modela experimentos cuyo resultado es incierto: una tirada de dado, una
carta extraída, un juego jugado. Este capítulo establece el vocabulario sobre finitos.
conjuntos, luego presenta el personaje central de todos \omterm{def:g10:proba:distribution}{probabilidad} teoría:
el \emph{variable aleatoria} y su \emph{\omterm{def:g11:prob:expectation}{esperanza}}. Estas nociones son
revisado y profundizado en \cref{ch:g12:randvar}y condicional
las \omterm{def:g10:proba:distribution}{probabilidades} siguen \cref{ch:g12:condprob}.

\section{Probabilidad en un conjunto finito}

\begin{definition}[Modelo de probabilidad]\label{def:g11:prob:model}
Un experimento con un número finito de resultados se modela por su \emph{muestra
espacio}\index{muestra
espacio} $\Omega = \{\omega_1, \dots, \omega_n\}$ (el
conjunto de resultados) junto con un \emph{\omterm{def:g10:proba:distribution}{probabilidad}} $\P $ asignando a cada
resultado $\omega_i $ un numero $ p_i \geq 0$, con $ p_1 + \dots + p_n = 1$. Un
\emph{evento}\index{evento} es un subconjunto $ A \subseteq \Omega $, y
\[
\P(A) = \sum_{\omega_i \in A} p_i .
\]
Cuando todos los resultados son igualmente probables ($ p_i = \frac1n $: el \emph{uniforme
modelo}\index{uniforme
modelo}),
\[
\P(A) = \frac{\text{number of outcomes in } A}
{\text{number of outcomes in } \Omega}.
\]
\end{definition}

\begin{example}\label{ex:g11:prob:twodice}
Tira dos dados distinguibles: $\Omega $ es el conjunto de los $36$ pares ordenados
$(i, j)$ con $1 \leq i, j \leq 6$, todos igualmente probables. El suceso
$ S $ = ``la suma es $7$'' contiene el $6$ pares
$(1,6)$, $(2,5)$, $(3,4)$, $(4,3)$, $(5,2)$, $(6,1)$, entonces
$\P(S) = \frac{6}{36} = \frac16$. El suceso ``la suma es $12$'' contiene
solo $(6,6)$: \omterm{def:g10:proba:distribution}{probabilidad} $\frac{1}{36}$. Eligiendo el \emph{bien}
\omterm{def:g10:proba:events}{espacio muestral} --- pares ordenados, no desordenados --- es lo que hace que el
modelo uniforme aplicable.
\end{example}

\begin{proposition}[Reglas de probabilidad]\label{prop:g11:prob:rules}
Para sucesos $ A, B \subseteq \Omega $:
\[
\P(\emptyset) = 0, \qquad \P(\Omega) = 1, \qquad
\P(\bar A) = 1 - \P(A),
\]
\[
\P(A \cup B) = \P(A) + \P(B) - \P(A \cap B),
\]
dónde $\bar A $ es el \emph{\omterm{def:g10:proba:operations}{complemento}} de $ A $ (todos los resultados que no están en $ A $).
En particular, $\P(A \cup B) = \P(A) + \P(B)$ cuando $ A $ y $ B $ son
\emph{\omterm{def:g10:proba:operations}{incompatibles}} ($ A \cap B = \emptyset $).
\end{proposition}

\begin{proof}
el vacio suceso no contiene ningún resultado (suma vacía), y $\Omega $ contiene
todos ellos (suma total $1$). Cada resultado está exactamente en uno de $ A $,
$\bar A $, entonces $\P(A) + \P(\bar A) = 1$. Para el \omterm{def:g10:numbers:interunion}{unión}: sumando el $ p_i $
encima $ A $ y luego encima $ B $ cuenta los resultados de $ A \cap B $ dos veces;
restando $\P(A \cap B)$ corrige la doble cuenta.
\end{proof}

\begin{example}
En una clase, $60\%$ estudiar música ($ M $), $45\%$ practicar un deporte ($ S $), y
$25\%$ haz ambas cosas. La proporción que hace al menos uno es
$\P(M \cup S) = 0.60 + 0.45 - 0.25 = 0.80$, entonces $20\%$ no hacer ninguna de las dos cosas
(\omterm{def:g10:proba:operations}{complemento}).
\end{example}

\section{Variables aleatorias}

\begin{definition}[Variable aleatoria, distribución]\label{def:g11:prob:rv}
A \emph{variables aleatorias}\index{variables aleatorias} en $\Omega $ es una \omterm{def:g11:func:function}{función}
$ X $ asignando un número a cada resultado. Es
\emph{distribución}\index{distribución} es la tabla de sus posibles
valores $ x_1, \dots, x_k $ junto con sus \omterm{def:g10:proba:distribution}{probabilidades}
\[
\P(X = x_i) = \P\bigl(\{\omega : X(\omega) = x_i\}\bigr),
\qquad
\sum_{i=1}^k \P(X = x_i) = 1 .
\]
\end{definition}

\begin{example}\label{ex:g11:prob:game}
Un juego: pagar $2$ euros, tira un dado justo y recibe el valor mostrado si
es al menos $5$, nada más. Sea $ G $ ser el \emph{ganar} (recibido
menos pagado). Los resultados $1, 2, 3, 4$ dar $ G = -2$; el resultado $5$
da $ G = 3$; el resultado $6$ da $ G = 4$. El \omterm{def:g11:prob:rv}{distribución} de $ G $:
\begin{center}
\begin{tabular}{c|ccc}
$ g $ & $-2$ & $3$ & $4$\\
\hline\rule{0pt}{11pt}%
$\P(G = g)$ & $\dfrac46$ & $\dfrac16$ & $\dfrac16$
\end{tabular}
\end{center}
\end{example}

\begin{method}[Construyendo una tabla de distribución]\label{met:g11:prob:table}
Enumere los posibles valores de $ X $; Para cada valor, recopile los resultados.
producirlo y sumar sus \omterm{def:g10:proba:distribution}{probabilidades}; comprobar que la fila de
sumas de \omterm{def:g10:proba:distribution}{probabilidades} a $1$.
\end{method}

\section{Esperanza}

\begin{definition}[Esperanza]\label{def:g11:prob:expectation}
El \emph{esperanza}\index{esperanza} (o \omterm{def:g11:stat:mean}{media}) de $ X $ es el promedio
de sus valores ponderados por sus \omterm{def:g10:proba:distribution}{probabilidades}:
\[
\E(X) = \sum_{i=1}^{k} \P(X = x_i)\; x_i .
\]
\end{definition}

\begin{remark}
$\E(X)$ es lo que la \omterm{def:g11:stat:mean}{media} de un gran número de repeticiones independientes de
el experimento sería --- la versión precisa de esta afirmación es la
ley de los grandes números, demostrada en \cref{ch:g12:sums}. Tenga en cuenta lo formal
analogía con la \omterm{def:g11:stat:mean}{media} de un conjunto de datos con \omterm{def:g10:stats:series}{frecuencias}
(\cref{def:g11:stat:mean}): las \omterm{def:g10:proba:distribution}{probabilidades} juegan el papel de relativo
\omterm{def:g10:stats:series}{frecuencias} $\frac{n_i}{n}$.
\end{remark}

\begin{example}\label{ex:g11:prob:gameexp}
Para el juego de \cref{ex:g11:prob:game}:
\[
\E(G) = \frac46 \times (-2) + \frac16 \times 3 + \frac16 \times 4
= \frac{-8 + 3 + 4}{6} = -\frac16 \approx -0.17 .
\]
En promedio, el jugador pierde aproximadamente $17$ centavos por juego: el juego es
desfavorable. un juego se llama \emph{justo} cuando la ganancia esperada es $0$.
\end{example}

\begin{proposition}[Linealidad]\label{prop:g11:prob:linearity}
Para todos los reales $ a, b $:
\[
\E(aX + b) = a\,\E(X) + b .
\]
\end{proposition}

\begin{proof}
la variable $ aX + b $ toma el valor $ a x_i + b $ con \omterm{def:g10:proba:distribution}{probabilidad}
$\P(X = x_i)$, entonces
\[
\E(aX + b) = \sum_i \P(X = x_i)(a x_i + b)
= a \sum_i \P(X = x_i)\,x_i + b \underbrace{\sum_i \P(X =
x_i)}_{=\,1} = a\,\E(X) + b .
\]
\end{proof}

\section{Varianza y desviación estándar}

\begin{definition}[Diferencia]\label{def:g11:prob:variance}
El \emph{diferencia}\index{diferencia} de $ X $ mide la propagación de su
valores alrededor del \omterm{def:g11:prob:expectation}{esperanza} $ m = \E(X)$:
\[
\V(X) = \E\bigl((X - m)^2\bigr)
= \sum_{i=1}^{k} \P(X = x_i)\,(x_i - m)^2,
\qquad
\sigma(X) = \sqrt{\V(X)} .
\]
\end{definition}

\begin{proposition}[Fórmula de atajo]\label{prop:g11:prob:konig}
\[
\V(X) = \E(X^2) - \E(X)^2 .
\]
\end{proposition}

\begin{proof}
Expandir $(x_i - m)^2 = x_i^2 - 2m\,x_i + m^2$ dentro de la suma:
\[
\V(X) = \sum_i \P(X = x_i)\,x_i^2
- 2m \sum_i \P(X = x_i)\,x_i + m^2\sum_i \P(X = x_i)
= \E(X^2) - 2m^2 + m^2 ,
\]
cual es $\E(X^2) - m^2$. Esta es la versión de variable aleatoria del
identidad probada para conjuntos de datos en \cref{prop:g11:stat:shortcut}.
\end{proof}

\begin{proposition}[Transformación afín]\label{prop:g11:prob:affine}
$\V(aX + b) = a^2\,\V(X)$ y $\sigma(aX + b) = \abs a\,\sigma(X)$.
\end{proposition}

\begin{proof}
Por linealidad, $ aX + b $ se desvía de su \omterm{def:g11:prob:expectation}{esperanza} $ a m + b $ por
$(aX + b) - (am + b) = a(X - m)$: cada desviación se escala por $ a $, por lo tanto
cada desviación al cuadrado por $ a^2$, y también lo es su peso \omterm{def:g11:stat:mean}{media}. el
turno $ b $ ha desaparecido --- cambiar una variable no cambia su
difundir.
\end{proof}

\begin{example}\label{ex:g11:prob:variance}
Para el juego de \cref{ex:g11:prob:game}:
$\E(G^2) = \frac46 \times 4 + \frac16 \times 9 + \frac16 \times 16
= \frac{16 + 9 + 16}{6} = \frac{41}{6}$, entonces
\[
\V(G) = \frac{41}{6} - \left(-\frac16\right)^2
= \frac{41}{6} - \frac{1}{36} = \frac{245}{36} \approx 6.8,
\qquad
\sigma(G) \approx 2.6 .
\]
la perdida de $0.17$ por juego en promedio viene con swings de tamaño típico
$2.6$ euros --- \omterm{def:g11:prob:expectation}{esperanza} sin desviación estándar dice solo la mitad
la historia.
\end{example}

\begin{omfigure}
\begin{tikzpicture}
\begin{axis}[
  width=8.6cm, height=5.2cm,
  ybar, bar width=14pt,
  xlabel={$ g $}, ylabel={$\P(G = g)$},
  xmin=-3.5, xmax=5.5, ymin=0, ymax=0.8,
  xtick={-2,3,4}, ytick={1/6, 4/6},
  yticklabels={$\tfrac16$,$\tfrac46$},
  tick label style={font=\small},
  label style={font=\small}]
  \addplot[fill=omDef!40, draw=omDef] coordinates {
    (-2,0.6667) (3,0.1667) (4,0.1667)};
  \draw[omThm, dashed, thick] (axis cs:-0.1667,0) --
    (axis cs:-0.1667,0.75);
  \node[omThm, font=\small, anchor=south west] at (axis cs:-0.05,0.62)
    {$\E(G)$};
\end{axis}
\end{tikzpicture}

{\small El \omterm{def:g11:prob:rv}{distribución} de la ganancia $ G $ de
\cref{ex:g11:prob:game}, y su \omterm{def:g11:prob:expectation}{esperanza} $\E(G) = -\frac16$ (discontinuo):
el punto de equilibrio del \omterm{def:g10:proba:distribution}{probabilidad} masas.}
\end{omfigure}

\section{Ejercicios}

\begin{exercise}[$\star $]\label{exo:g11:prob:1}
Se extrae una carta de un estándar. $52$-baraja de cartas. Calcular el \omterm{def:g10:proba:distribution}{probabilidad}
del dibujo: un corazón; un rey; un corazón o un rey; ni un corazón ni un
rey.
\end{exercise}

\begin{exercise}[$\star $]\label{exo:g11:prob:2}
Se lanzan dos dados justos y su suma $ S $ está grabado. Usando el
$36$-modelo de resultados de \cref{ex:g11:prob:twodice}, calcular $\P(S = 6)$,
$\P(S \geq 10)$ y $\P(S \text{ is even})$.
\end{exercise}

\begin{exercise}[$\star $]\label{exo:g11:prob:3}
A variable aleatoria tiene \omterm{def:g11:prob:rv}{distribución}
\begin{center}
\begin{tabular}{c|cccc}
$ x $ & $-1$ & $0$ & $2$ & $5$\\
\hline
$\P(X=x)$ & $0.3$ & $0.2$ & $0.4$ & $ p $
\end{tabular}
\end{center}
Encontrar $ p $, luego calcular $\E(X)$.
\end{exercise}

\begin{exercise}[$\star $]\label{exo:g11:prob:4}
$ X $ satisface $\E(X) = 3$ y $\V(X) = 4$. Dale la \omterm{def:g11:prob:expectation}{esperanza}, varianza
y desviación estándar de $ Y = 2X - 5$.
\end{exercise}

\begin{exercise}[$\star\star $]\label{exo:g11:prob:5}
Una rueda de la fortuna tiene $10$ sectores iguales: $5$ espectáculo $0$ euros, $3$ espectáculo
$2$ euros, $2$ espectáculo $5$ euros. Costos de juego $2$ euros. Dale el
\omterm{def:g11:prob:rv}{distribución} de la ganancia $ G $ (ganancias menos apuesta), calcular $\E(G)$, y
decidir si el juego es favorable.
\end{exercise}

\begin{exercise}[$\star\star $]\label{exo:g11:prob:6}
Una urna contiene $3$ rojo y $2$ bolas azules; se extraen dos bolas
\emph{sin} reemplazo. Sea $ X $ Sea el número de bolas rojas extraídas.
Construye el \omterm{def:g11:prob:rv}{distribución} de $ X $ (enumere los $20$ sorteos ordenados, o utilizar un
\omterm{met:g10:proba:tree}{árbol}), y calcular $\E(X)$.
\end{exercise}

\begin{exercise}[$\star\star $]\label{exo:g11:prob:7}
para la rueda de \cref{exo:g11:prob:5}, calcular $\V(G)$ y $\sigma(G)$.
\end{exercise}

\begin{exercise}[$\star\star $]\label{exo:g11:prob:8}
Una compañía de seguros vende una póliza por $80$ euros. Con \omterm{def:g10:proba:distribution}{probabilidad}
$0.05$ el cliente reclama $1000$ euros, con \omterm{def:g10:proba:distribution}{probabilidad} $0.01$ el cliente
reclamaciones $5000$ euros, por lo demás nada. Sea $ B $ ser el beneficio de la empresa
en una sola política. Dale el \omterm{def:g11:prob:rv}{distribución} de $ B $, calcular $\E(B)$, y
interpretar.
\end{exercise}

\begin{exercise}[$\star\star $]\label{exo:g11:prob:9}
Una pregunta de opción múltiple tiene $4$ respuestas, una correcta. una respuesta correcta
puntuaciones $+1$; para desalentar las adivinanzas, una respuesta incorrecta puntúa $ x < 0$.
Un estudiante responde al azar. ¿Para qué valor de $ x $ es del estudiante
puntuación esperada cero?
\end{exercise}

\begin{exercise}[$\star\star $]\label{exo:g11:prob:10}
Se lanzan dos monedas justas. Ann propone: ``si ambas monedas coinciden, te pago
$1$ euro; si difieren me pagas $1$ euro''. Calcule el valor esperado
ganancia de cada jugador. ¿Es el juego justo? Misma pregunta con tres monedas,
donde ann paga $2$ euros si coinciden los tres.
\end{exercise}

\begin{exercise}[$\star\star\star $]\label{exo:g11:prob:11}
una loteria vende $1000$ entradas en $5$ euros cada uno. Un boleto gana
$2000$ euros, ganan cinco billetes $200$ euros, ganan veinte billetes $25$
euros.
\begin{enumerate}
  \item Dale el \omterm{def:g11:prob:rv}{distribución} de los ingresos totales del organizador menos
        premios, y de la ganancia de un solo jugador $ G $.
  \item Calcular $\E(G)$ y el beneficio total esperado del organizador.
        Comprueba que las dos respuestas sean consistentes.
\end{enumerate}
\end{exercise}

\section{Problema: el juego inacabado}

\begin{problem}[{Problema del fin de semana --- Pascal, Fermat y la feria
Escisión de una apuesta interrumpida: el nacimiento de la esperanza, la magia
de linealidad, y lo que la esperanza no puede ver}]
\label{pb:g11:prob:1}
En 1654 el Caballero de M.éjemé --- el mismo jugador cuyos dados
Las apuestas abrieron el volumen de la escuela secundaria. \omterm{def:g10:proba:distribution}{probabilidad} problema ---
le hizo a Blaise Pascal una pregunta más profunda: dos personas igualmente hábiles
Los jugadores deben detener su serie con el puntaje de $5$ a $3$, seis
gana tomando el $64$-olla de pistola. \emph{¿Cómo debe ser la olla?
dividir de manera justa?} Las cartas entre Pascal y Fermat respondiendo
fundó el concepto central de este capítulo: \omterm{def:g11:prob:expectation}{esperanza}
(\cref{def:g11:prob:expectation}). Volverás a derivar su
responde en ambos sentidos, luego conoce el truco de fiesta de la linealidad y
esperanzaHay un punto ciego.

\medskip
\noindent\textbf{Part I --- \omterm{def:g11:prob:expectation}{esperanza} fluency.}
\begin{enumerate}
  \item Tira un dado justo; tu ganancia es (cara $- 3$) euros. construir
        el \omterm{def:g11:prob:rv}{distribución} tabla de ganancia $ G $
        (\cref{met:g11:prob:table}) y calcular $\E(G)$.
  \item Calcular $ V(G)$ con la fórmula abreviada
        (\cref{prop:g11:prob:konig}) y $\sigma(G)$ (dos
        decimales).
  \item Sin ninguna tabla nueva, dar $\E(2G - 1)$ y
        $ V(2G - 1)$ (\cref{prop:g11:prob:affine}).
  \item Un puesto cobra $ m $ euros para tirar dos dados y casillas
        te la suma en euros. cual es el \emph{justo} precio
        $ m $? El puesto cobra $8$: ¿Cuál es su pérdida esperada?
        por juego?
  \item Un portátil que vale la pena $800$ euros tiene un $5\,\%$ posibilidad de
        destrucción cada año. Calcular el seguro justo
        prima y el beneficio esperado del asegurador por póliza
        si se carga $60$ euros.
\end{enumerate}

\medskip
\noindent\textbf{Part II --- Splitting the pot.}
Puntaje $5$--$3$, primero en $6$ gana el $64$ pistolas, balas son
lanzamientos de moneda justos.
\begin{enumerate}[resume]
  \item Dos tentadoras respuestas erróneas: dividir $5 : 3$
        (proporcionalmente a las rondas ganadas), o dar todo a
        el líder. Di en una frase cada uno lo que es injusto.
        sobre ellos.
  \item El método de Fermat: el jugador A necesita $1$ más victoria, B necesita
        $3$; imagina jugar exactamente $3$ más rondas independientemente.
        Enumere el $8$ resultados igualmente probables, marque quién gana
        serie en cada uno, y deducir la división justa de la $64$
        pistolas.
  \item El punto sutil que Fermat tuvo que defender: el juego real
        podría detenerse después de una ronda --- ¿por qué sigue siendo así?
        legítimo para razonar sobre todos $3$ rondas imaginarias?
  \item El método de Pascal, al revés desde el final: calcular A
        participación justa en las puntuaciones (hipotéticas) $5$--$5$, entonces
        $5$--$4$, entonces $5$--$3$, cada vez promediando los dos
        futuros igualmente probables. Compárese con la pregunta 7.
  \item Misma pregunta en la partitura. $4$--$3$: cuantos imaginarios
        rondas, qué resultados le dan a B la serie y cuál es
        la división justa?
  \item Establezca el principio que comparten los dos métodos: la feria.
        El valor de un futuro incierto es su \emph{\omterm{def:g11:prob:expectation}{esperanza}}
        --- y nombre dos industrias modernas que funcionan enteramente con
        (la pregunta 5 \omterm{def:g10:proba:sample}{muestra} uno).
  \item Un control de cordura para cada puntuación: explique, con el
        linealidad de \omterm{def:g11:prob:expectation}{esperanza}
        (\cref{prop:g11:prob:linearity}), por qué los dos jugadores
        las partes justas siempre deben sumar exactamente el
        $64$-olla de pistola.
\end{enumerate}

\medskip
\noindent\textbf{Part III --- Linearity's party trick.}
\begin{enumerate}[resume]
  \item Recuperar $\E(\text{sum of two dice}) = 7$ en una linea
        con linealidad (\cref{prop:g11:prob:linearity}), No
        $36$-mesa de celdas. Para el \emph{varianza} de la suma,
        la regla $ V(X + Y) = V(X) + V(Y)$ necesita que los dados sean
        independiente (admitido): computarlo.
  \item La fiesta del hat-check: $ n $ los invitados arrojan sus sombreros en un
        apila y cada uno toma uno al azar. Sea $ X $ contar
        los invitados que recuperan su propio sombrero. por uno fijo
        invitado, ¿cuál es el \omterm{def:g10:proba:distribution}{probabilidad} de recuperar los suyos
        sombrero? Sumando estos $ n $ posibilidades (¡linealidad!), calcular
        $\E(X)$. La respuesta debería sorprenderte: no
        depender de $ n $.
  \item Verificar $\E(X) = 1$ por fuerza bruta para $ n = 2$ y
        $ n = 3$ (enumere el $2$ y $6$ igualmente probable
        asignaciones de sombreros y promediar los conteos).
  \item Las recuperaciones de los invitados están lejos de ser independientes (si
        $ n - 1$ los invitados recuperaron sus propios sombreros, el último
        también debe hacerlo). ¿Dónde exactamente se realizó el cálculo de la pregunta 14?
        \emph{no} ¿Requiere independencia? Linealidad del estado
        superpotencia.
\end{enumerate}

\medskip
\noindent\textbf{Part IV --- What \omterm{def:g11:prob:expectation}{esperanza} cannot see.}
\begin{enumerate}[resume]
  \item Dos $5$-Tarjetas rasca y gana en euros: la tarjeta A gana $10$ euros con
        \omterm{def:g10:proba:distribution}{probabilidad} $\frac12$; la tarjeta B gana $1\,000$ euros con
        \omterm{def:g10:proba:distribution}{probabilidad} $0.005$. Compruebe que ambas ganancias tengan la
        lo mismo \omterm{def:g11:prob:expectation}{esperanza}, luego calcula ambos variaciones. cual
        ¿Qué compran las personas reales y qué compran?
  \item El juego de San Petersburgo: voltear hasta que salga la primera cruz;
        si la primera cruz ocurre al girar $ n $ tu recibes
        $2^n $ euros. Calcular la contribución de cada $ n $ hacia \omterm{def:g11:prob:expectation}{esperanza} y la \omterm{def:g11:prob:expectation}{esperanza} sí mismo. ¿Quieres
        pagar $1\,000$ euros para jugar una vez? ¿Qué significa el choque?
        enseñar sobre \omterm{def:g11:prob:expectation}{esperanza}?
  \item Seguro por parte del cliente: pagar $20$ euros, o
        enfrentar un $1\,\%$ riesgo de perder $1\,000$? comparar los dos
        \omterm{def:g11:prob:expectation}{esperanzas} de heredary explique por qué comprar el seguro puede
        seguir siendo racional (¿qué no puede hacer una sola familia?)
        ``promedio superior''?).
  \item Final --- esperanzabiografía de, una frase por
        capítulo: nacido en 1654 para dividir una olla; su superpotencia,
        linealidad, que suma \omterm{def:g11:prob:expectation}{esperanzas} de heredar sin preguntar
        independencia (los sombreros); su punto ciego, el riesgo, que
        varianza medidas (tarjetas, San Petersburgo, seguros);
        y su coronación el próximo año, cuando la ley de los grandes
        Los números explican exactamente por qué la casa siempre gana.
\end{enumerate}
\end{problem}
